\documentclass[12pt]{article}
\usepackage{amsmath,amssymb,amsfonts}
\usepackage[margin=1in]{geometry}
\begin{document}
\begin{center}
{\large\bfseries 35th Irish Mathematical Olympiad\\7 May 2022}
\end{center}
\bigskip\noindent\textbf{Paper 1}\bigskip
\begin{enumerate}
\item For $n$ a positive integer, $n! = 1 \cdot 2 \cdot 3 \cdots (n-1) \cdot n$ is the product of the positive integers from 1 to $n$.

Determine, with proof, all positive integers $n$ for which $n! + 3$ is a power of 3.

\item Let $ABCD$ be a square and let $\Gamma$ denote the circle with diameter $CD$. A tangent line is drawn to the circle $\Gamma$ from $B$, meeting the circle $\Gamma$ at $E$ and intersecting the segment $AD$ at $K$.

Prove that $|AD| = 4|KD|$.

\item Let $n \ge 3$ be an integer and let $(p_1, p_2, p_3, \ldots, p_n)$ be a permutation of $\{1, 2, 3, \ldots, n\}$. For this permutation we say that $p_t$ is a turning point if $2 \le t \le n - 1$ and
\[
(p_t - p_{t-1})(p_t - p_{t+1}) > 0.
\]
For example, for $n = 8$, the permutation $(2, 4, 6, 7, 5, 1, 3, 8)$ has two turning points: $p_4 = 7$ and $p_6 = 1$.

For fixed $n$, let $q(n)$ denote the number of permutations of $\{1, 2, 3, \ldots, n\}$ with exactly one turning point. Find all $n \ge 3$ for which $q(n)$ is a perfect square.

\item Let $\mathbb{N}$ denote the strictly positive integers. A function $f : \mathbb{N} \mapsto \mathbb{N}$ has the following properties which hold for all $n \in \mathbb{N}$:
\begin{enumerate}
\item[(a)] $f(n) < f(n+1)$;
\item[(b)] $f(f(f(n))) = 4n$.
\end{enumerate}
Find $f(2022)$.

\item Let $\triangle ABC$ be a triangle with circumcentre $O$. The perpendicular line from $O$ to $BC$ intersects line $BC$ at $M$ and line $AC$ at $P$, and the perpendicular line from $O$ to $AC$ intersects line $AC$ at $N$ and line $BC$ at $Q$. Let $D$ be the intersection point of lines $PQ$ and $MN$. Construct the parallelogram $PCQJ$ with $PJ \parallel CQ$ and $QJ \parallel CP$.

Prove the following:
\begin{enumerate}
\item[(a)] The points $A$, $B$, $O$, $P$, $Q$, $J$ are all on the same circle.
\item[(b)] Line $OD$ is perpendicular to line $CJ$.
\end{enumerate}
\end{enumerate}

\newpage
\begin{center}
{\large\bfseries 35th Irish Mathematical Olympiad\\7 May 2022}
\end{center}
\bigskip\noindent\textbf{Paper 2}\bigskip
\begin{enumerate}
\setcounter{enumi}{5}
\item Suppose $a, b, c$ are real numbers such that $a + b + c = 1$. Prove that
\[
a^3 + b^3 + c^3 + 3(1-a)(1-b)(1-c) = 1.
\]

\item The four vertices of quadrilateral $ABCD$ lie on the circle with diameter $AB$. The diagonals of $ABCD$ intersect at $E$, and the lines $AD$ and $BC$ intersect at $F$. Line $FE$ meets $AB$ at $K$ and line $DK$ meets the circle again at $L$.

Prove that $CL$ is perpendicular to $AB$.

\item The equation $AB \times CD = EFGH$, where each of the letters $A$, $B$, $C$, $D$, $E$, $F$, $G$, $H$ represents a different digit and the values of $A$, $C$ and $E$ are all nonzero, has many solutions, e.g., $46 \times 85 = 3910$. Find the smallest value of the four-digit number $EFGH$ for which there is a solution.

\item Let $k$ be a positive integer and let $x_0, x_1, x_2, \ldots$ be an infinite sequence defined by the relationship
\[
x_0 = 0, \quad x_1 = 1, \quad x_{n+1} = kx_n + x_{n-1} \text{ for all } n \ge 1.
\]
\begin{enumerate}
\item[(a)] For the special case $k = 1$, prove that $x_{n-1}x_{n+1}$ is never a perfect square for $n \ge 2$.
\item[(b)] For the general case of integers $k \ge 1$, prove that $x_{n-1}x_{n+1}$ is never a perfect square for $n \ge 2$.
\end{enumerate}

\item Let $n \ge 5$ be an odd number and let $r$ be an integer such that $1 \le r \le (n-1)/2$. In a sports tournament, $n$ players take part in a series of contests. In each contest, $2r+1$ players participate, and the scores obtained by the players are the numbers
\[
-r, -(r-1), \ldots, -1, 0, 1, \ldots, r-1, r
\]
in some order. Each possible subset of $2r+1$ players takes part together in exactly one contest. Let the final score of player $i$ be $S_i$, for each $i = 1, 2, \ldots, n$. Define $N$ to be the smallest difference between the final scores of two players, i.e.,
\[
N = \min_{i < j} |S_i - Sj|.
\]
Determine, with proof, the maximum possible value of $N$.
\end{enumerate}
\end{document}
